%% notes-tex/database/yeastract-publications/sections/1-file.tex
%% [[experiments.037-yeastract.publications]]
%% https://github.com/Mjvolk3/torchcell/tree/main/notes-tex/database/yeastract-publications/sections/1-file.tex
%% SOURCE: numbers and tables from
%% experiments/037-yeastract/scripts/yeastract_publication_tables.py; prose hand-authored

\section{What the file holds\secstatus{todo}}\label{sec:file}

The source is the \org{S. cerevisiae} member of the YEASTRACT+ 2022 regulation
archive, which holds one flat file for each of eleven yeast species:

\begin{quote}\small
archive: \file{YeastractPlus_2022_all_regs.zip}\\
sha256: \texttt{2f2e9f8b3c905ae333814f04161d0ffd\allowbreak 2fabf00cc9878c1bd86a8f5e366ebb3e}\\
member: \file{YeastractPlus_2022_Scerevisiae_regs_flatfile.tsv.gz}
\end{quote}

\noindent
The member is a tab-separated table of \ytNRows{} rows and eleven columns with no header row,
and no row is repeated. The column meanings used here are read off the values:
the systematic and standard names of a transcription factor, the systematic and
standard names of a target gene, a PMID, a sign (\texttt{Positive},
\texttt{Negative} or \texttt{N/A}), an evidence label (\texttt{Direct},
\texttt{Indirect} or \texttt{N/A}), an assay, a condition group, a condition
subgroup and a free-text condition.

A row is one piece of evidence, and one factor-target pair recurs once for each
paper, assay and condition that reports it. The \ytNRows{} rows name \ytNTfs{}
transcription factors and \ytNTargets{} targets in \ytNPairs{} distinct pairs.
\ytNIndirectRows{} rows are labeled \texttt{Indirect}, \ytNDirectRows{} are
labeled \texttt{Direct} and \ytNEvidenceNaRows{} carry no label
(Table~\ref{tab:summary}).

%% GENERATED FILE -- do not hand-edit.
%% SOURCE: experiments/037-yeastract/scripts/yeastract_publication_tables.py
\begin{table}[H]\centering
\footnotesize
\caption[]{Papers and rows in the YEASTRACT+ flat file by relation to the SPELL
archive. \emph{Rows} counts flat-file rows credited to the papers and \emph{\%}
is their share of all rows. \emph{Direct} and \emph{indirect} count rows by the
evidence label in column 7 of the file; rows labeled \texttt{N/A} are in
neither. \emph{DOI} and \emph{PMC id} count the papers whose PubMed record
carries one. \emph{Years} is the range of publication years.}\label{tab:summary}
\begin{tabular}{lrrrrrrrl}
\toprule
& papers & rows & \% & direct & indirect & DOI & PMC id & years \\
\midrule
same PMID in SPELL & 84 & 75,342 & 21.2 & 3,493 & 71,846 & 84 & 55 & 1997--2019 \\
reanalysis of a SPELL study & 1 & 66,920 & 18.8 & 0 & 66,920 & 1 & 1 & 2010--2010 \\
not in SPELL & 1,586 & 213,919 & 60.1 & 87,662 & 126,006 & 1,547 & 920 & 1976--2022 \\
\midrule
\textbf{all} & 1,671 & 356,181 & 100.0 & 91,155 & 264,772 & 1,632 & 976 & 1976--2022 \\
\bottomrule
\end{tabular}
\end{table}


\subsection{Where each field comes from\secstatus{todo}}\label{sec:fields}

The PMID is column 5 of the file, and the row, factor, pair, evidence and
condition counts of a paper are taken over the rows that carry its PMID. The
file gives no citation, so the authors, title, journal, DOI and PMC id come
from the PubMed record for that PMID, retrieved through the NCBI E-utilities
\texttt{esummary} endpoint and stored beside the manifest in
\file{experiments/037-yeastract/results/yeastract_pubmed_esummary.json}.

PubMed returns no record for \ytNNoPubmed{} of the \ytNPapers{} PMIDs, and that
row is listed by its id alone. \ytNDoi{} records carry a DOI and \ytNNoDoi{} do
not; those are linked by their PubMed record. \ytNPmc{} records carry a PMC id,
which means PubMed Central holds a full text. Whether that copy is in the
open-access subset served by the PMC OA API was not checked.

For \ytNNoAssayPapers{} papers the most frequent value of the assay column is
empty, and those papers show a dash in the assay column of Table~\ref{tab:all}.
